\section{Пошаговый алгоритм развёртывания}
\headingtextsep

Пошаговый алгоритм развёртывания должен обеспечивать предсказуемый запуск всех компонентов системы и контроль ключевых точек готовности. Для системы «Контур» такой алгоритм строится вокруг контейнерного запуска через Docker Compose и включает подготовку окружения, сборку образов, запуск базы данных, старт прикладных сервисов и проверку доступности интерфейсов. UML-схема алгоритма приведена в приложении А на рисунке \ref{fig:appendix-deployment-algorithm}.

\textheadingsep
\subsection{Описание алгоритма}
\headingtextsep

Алгоритм развёртывания включает следующие основные шаги:

\begin{enumerate}
\item проверка наличия Docker Engine или Docker Desktop вместе с Docker Compose;
\item получение исходного кода проекта и подготовка конфигурации окружения на основе файла .env.example;
\item проверка доступности требуемых сетевых портов и готовности локальной машины к запуску контейнеров;
\item сборка образов для API-сервиса, web-клиента и MCP-моста;
\item запуск контейнера PostgreSQL и ожидание успешного прохождения healthcheck;
\item запуск API-сервиса после готовности базы данных;
\item запуск web-клиента и MCP-моста после старта API;
\item проверка доступности web-интерфейса и HTTP API;
\item анализ журналов и повторный запуск проблемного компонента в случае ошибки;
\item перевод системы в рабочее состояние после успешной проверки всех сервисов.
\end{enumerate}

\textheadingsep
\subsection{Логика переходов между шагами}
\headingtextsep

В алгоритме присутствуют контрольные точки, оформленные как условия перехода. Если на начальном этапе отсутствуют Docker и Docker Compose, развёртывание не может быть продолжено до установки соответствующих средств. Если заняты необходимые сетевые порты, конфигурация машины должна быть исправлена до запуска контейнеров. Отдельная проверка выполняется после старта PostgreSQL: пока база данных не прошла healthcheck, запуск зависящих сервисов откладывается.

На завершающем этапе оценивается доступность пользовательского интерфейса и API. Если хотя бы один из сервисов недоступен, алгоритм переходит к диагностике журналов и повторному запуску отказавшего компонента. Если все проверки проходят успешно, развёртывание считается завершенным и система переводится в рабочий режим.
